%=============================================================================!
% 14.1  FORCE ON A WALL                                                       !
%=============================================================================!
\section{Force on a wall}
\label{sec:pr-wall}

A bicycle pump with its outlet closed pushes back on the hand: the air
inside presses on the piston, and the harder the air is squeezed the
harder it presses. The push divided by the area of the piston is the
pressure of the air. This section finds where that push comes from in a
gas of atoms, and measures it.

\subsection*{Pressure and its units}

The pressure on a flat wall of area $A$ that feels the force
$F_\perp$ at right angles to it is
\begin{equation*}
   P = \frac{F_\perp}{A} ,
\end{equation*}
in newtons per square metre, or pascals. The air around us presses with
about $10^5$\,\si{\pascal}, one bar; deep in the Earth and in squeezed
solids pressures reach gigapascals, \SI{1}{\giga\pascal} being $10^9$
pascals, or $10^4$ bar. In the units of this book a pressure is an
energy per volume, \si{\electronvolt\per\angstrom\cubed}, since a force is
an energy per length and an area a length squared; one
\si{\electronvolt\per\angstrom\cubed} is \SI{160.218}{\giga\pascal}
(\cref{ex:pr-units}).

\subsection*{The push of bouncing atoms}

\begin{figure}[htb]
\centering
\begin{tikzpicture}[line cap=round, line join=round, scale=0.9]
   \fill[housegrey!15] (6.0,-0.2) rectangle (6.4,2.2);
   \draw[line width=0.6pt] (0,-0.2) -- (0,2.2);
   \draw[line width=0.6pt] (6.0,-0.2) -- (6.0,2.2);
   \node[font=\small, anchor=north] at (0,-0.25) {$x = 0$};
   \node[font=\small, anchor=north] at (6.0,-0.25) {$x = L$};
   \fill[accent] (5.55,1.45) circle (0.15);
   \draw[-{Stealth[length=5pt]}, line width=0.8pt, accent] (4.2,1.45) --
      node[above, font=\small] {$v_x$} (5.3,1.45);
   \fill[accent!45] (5.55,0.55) circle (0.15);
   \draw[-{Stealth[length=5pt]}, line width=0.8pt, accent!45] (5.3,0.55) --
      node[below, font=\small] {$-v_x$} (4.2,0.55);
   \node[font=\small, align=center] at (2.2,1.0)
      {$2mv_x$ to the wall\\once every $2L/v_x$};
\end{tikzpicture}
\caption{An atom bouncing between walls at $x = 0$ and $x = L$. Each
bounce off the right-hand wall turns $v_x$ into $-v_x$, and the atom
returns to that wall after travelling $2L$ along $x$.}
\label{fig:pr-bounce}
\end{figure}

One atom of mass $m$ moves in a box between walls at $x = 0$ and $x =
L$, its velocity along $x$ being $v_x$ (\cref{fig:pr-bounce}). A hard
wall takes no energy from the atom, so when it strikes the wall at $x =
L$ it comes back at the same speed with $v_x$ turned to $-v_x$: its
momentum along $x$ changes by $-2mv_x$, and by Newton's third law
(\cref{sec:ne-third}) the wall receives $+2mv_x$. The atom then crosses
the box and back, a distance $2L$ at the speed $v_x$, before it strikes
the same wall again; bounces off the walls normal to $y$ and $z$ turn
$v_y$ or $v_z$ and leave $v_x$ as it was. The bounces therefore come once
every $2L/v_x$. The wall feels a short kick at each bounce and nothing
between, and a force is the momentum delivered per unit time
(\cref{sec:ne-second}), so the force on the wall, averaged over many
bounces, is
\begin{equation*}
   \overline{F_\perp} = \frac{2mv_x}{2L/v_x} = \frac{mv_x^2}{L} .
\end{equation*}
For $N$ atoms that do not meet one another, each contributes its own
share, and the pressure on the wall, of area $A$, is
\begin{equation*}
   P = \frac{1}{A}\sum_i\frac{m_iv_{ix}^2}{L} = \frac{1}{V}\sum_i
   m_iv_{ix}^2 , \qquad V = AL .
\end{equation*}
Averaged over the velocities, each $\frac12m_iv_{ix}^2$ is $\frac12\kB T$
by equipartition (\cref{sec:sm-equipartition}), so each $m_iv_{ix}^2$ is
$\kB T$, the $N$ of them add to $N\kB T$, and
\begin{equation}
   PV = N\kB T ,
   \label{eq:pr-ideal}
\end{equation}
the ideal gas law that \cref{sec:sm-other} found by counting states. The
pressure of a gas is the momentum its atoms deliver to the walls per unit
time and area. The argument holds for molecules too, with $m$ the mass of
a molecule and $v_x$ the velocity of its centre of mass
(\cref{sec:ne-bodies}). At \SI{300}{\kelvin}, where $\kB T =
\SI{0.025852}{\electronvolt}$, and one bar,
\SI{6.2415e-7}{\electronvolt\per\angstrom\cubed} (\cref{ex:pr-units}),
\cref{eq:pr-ideal} gives $N/V = P/\kB T = \SI{2.414e-5}{\per\angstrom
\cubed}$: 24.1 molecules of air in a cube \SI{10}{\nano\metre}, or
\SI{100}{\angstrom}, on a side.

\subsection*{Measured}

A simulation has no hard walls, but soft ones serve: a spring that
pushes back an atom which has gone a distance $\delta$ beyond the wall
with the force $k\delta$, from the energy $\frac12k\delta^2$
(\texttt{virial.harmonic\_walls}). The force the atoms exert on each
wall is then the sum of these $k\delta$, recorded at every step. For 200
argon atoms that do not interact, in a cube of side \SI{40}{\angstrom}
with walls of stiffness \SI{5}{\electronvolt\per\angstrom\squared},
followed at fixed energy for \SI{100}{\pico\second} with steps of
\SI{5}{\femto\second}, the push on the six walls per unit area averages
\SI{53.47}{bar} (\cref{fig:pr-wall}(a)), and the two halves of the run
give 53.27 and \SI{53.67}{bar}. The kinetic energy of the run gives the
temperature \SI{123.5}{\kelvin} by \cref{eq:sm-kinetic-temperature}, with
$N_{\mathrm f} = 3N$ since the walls leave the total momentum free, and
\cref{eq:pr-ideal} gives
\SI{53.29}{bar}, within the difference between the halves.

\begin{figure}[htb]
\centering
\includegraphics{ch14_pressure/wall}
\caption{Pressure as the push on the walls. (a) 200 argon atoms that do
not interact, in a cube of \SI{40}{\angstrom} between soft walls: the push
per unit area on the six walls, averaged over each picosecond (thin) and
over the run so far (thick), against $N\kB T/V$ (dashed). (b) 256 argon
atoms that interact, in a cube of \SI{23.26}{\angstrom} near
\SI{135}{\kelvin}: the running means of the push per unit area and of the
pressure from the virial, \cref{eq:pr-virial}, against $N\kB T/V$
(dashed).}
\label{fig:pr-wall}
\end{figure}

Atoms that interact push differently. The 256 atoms of the liquid argon
of Chapter~13, held between walls of the same stiffness in a cube of
\SI{23.26}{\angstrom} under CSVR at \SI{135}{\kelvin}, push on the walls
with \SI{0.0624}{\giga\pascal} averaged over \SI{100}{\pico\second} (0.0616
and \SI{0.0633}{\giga\pascal} in the two halves). Their kinetic energy
gives the temperature \SI{133.3}{\kelvin} over the run, a little below the
\SI{135}{\kelvin} the thermostat holds on average, and with it
$N\kB T/V$ is only \SI{0.0374}{\giga\pascal} (\cref{fig:pr-wall}(b)). The atoms
repel one another at close range and the repulsion adds to the push on
the walls. The next section finds by how much.

\begin{keyidea}
Pressure is force per unit area. In a gas it is the momentum the atoms
deliver to the walls per unit time and area, $\sum_im_iv_{ix}^2/V$,
which averages to $N\kB T/V$ for atoms that do not interact. One
\si{\electronvolt\per\angstrom\cubed} is \SI{160.218}{\giga\pascal}, and
\SI{1}{\giga\pascal} is $10^4$ bar.
\end{keyidea}

\notebook{14}{let atoms bounce between walls and read their push}

\begin{selfcheck}
What pressure, in bar, do $10^4$ atoms that do not interact exert at
\SI{300}{\kelvin} in a cube of side \SI{100}{\angstrom}?
\end{selfcheck}
